\documentclass[12pt,a4paper]{article}

\usepackage[margin=2.5cm]{geometry}
\usepackage{hyperref}
\usepackage{graphicx}
\usepackage{listings}
\usepackage{xcolor}
\usepackage{booktabs}
\usepackage{enumitem}

\definecolor{codebg}{rgb}{0.95,0.95,0.95}
\lstset{
    backgroundcolor=\color{codebg},
    basicstyle=\ttfamily\small,
    breaklines=true,
    frame=single,
    numbers=left,
    numberstyle=\tiny\color{gray},
    keywordstyle=\color{blue},
    commentstyle=\color{green!50!black},
    stringstyle=\color{red!60!black},
}

\title{\textbf{Mini Game Hub} \\ \large Project Report}
\author{Shresth \and Aryan}
\date{\today}

\begin{document}

\maketitle
\tableofcontents
\newpage

% ─────────────────────────────────────────────────────────────────────
\section{External Libraries Used}

\begin{table}[h]
\centering
\begin{tabular}{ll}
\toprule
\textbf{Library} & \textbf{Purpose} \\
\midrule
\texttt{pygame-ce}     & Graphical user interface for all games and menus \\
\texttt{numpy}         & Board representation as 2-D arrays; vectorised win detection \\
\texttt{matplotlib}    & Post-game bar and pie chart visualisations \\
\texttt{pathlib}       & Cross-platform file path handling \\
\texttt{csv}           & Reading/writing \texttt{history.csv} \\
\texttt{subprocess}    & Calling \texttt{leaderboard.sh} from Python \\
\texttt{datetime}      & Timestamping game results \\
\texttt{collections.Counter} & Counting wins and game frequencies for charts \\
\bottomrule
\end{tabular}
\caption{Python libraries and their roles.}
\end{table}

Bash commands used: \texttt{sha256sum} / \texttt{shasum} (for SHA-256 hashing),
\texttt{awk} (text processing in \texttt{leaderboard.sh}), and standard utilities
(\texttt{sort}, \texttt{wc}, \texttt{tail}).

% ─────────────────────────────────────────────────────────────────────
\section{Features Implemented}

\begin{enumerate}[leftmargin=*]
    \item \textbf{Bash Authentication} (\texttt{main.sh}) -- Login and registration
          for two players with SHA-256 password hashing.  Stored in
          \texttt{users.tsv}.
    \item \textbf{Game Menu} (Pygame GUI) -- Interactive menu to choose
          Tic-Tac-Toe, Othello, or Connect Four.
    \item \textbf{Tic-Tac-Toe} -- 10$\times$10 board, 5-in-a-row to win.
          Win detection via NumPy sliding windows.
    \item \textbf{Othello (Reversi)} -- 8$\times$8 board with full move
          validation, disc flipping, turn skipping, and end-game scoring.
    \item \textbf{Connect Four} -- 7$\times$7 vertical grid with gravity,
          coin-drop animation, and NumPy-based win detection.
    \item \textbf{Result Recording} -- Each game appends winner, loser,
          date, and game name to \texttt{history.csv}.
    \item \textbf{Leaderboard} (\texttt{leaderboard.sh}) -- Formatted table
          showing per-user per-game wins, losses, and W/L ratio, sortable
          by any metric.
    \item \textbf{Matplotlib Visualisations} -- Top-5 players bar chart and
          most-played games pie chart displayed after each game.
    \item \textbf{Post-Game Loop} -- Players can choose to play again or quit.
\end{enumerate}

% ─────────────────────────────────────────────────────────────────────
\section{Code Logic Explanation}

\subsection{Generic Base Class (\texttt{base\_game.py})}
\texttt{BoardGame} is an abstract class providing:
\begin{itemize}
    \item A NumPy array (\texttt{self.board}) initialised to zeros.
    \item Turn tracking via \texttt{current\_player} (1 or 2).
    \item \texttt{switch\_turn()} toggles between players.
    \item \texttt{check\_win()} is abstract -- each game implements its own logic.
\end{itemize}

\subsection{Tic-Tac-Toe (\texttt{tictactoe.py})}
Player~1 places $+1$, Player~2 places $-1$.  Win detection uses
\texttt{numpy.lib.stride\_tricks.sliding\_window\_view} to extract all
$5$-cell windows horizontally and vertically, then sums each window.
A sum of $+5$ means Player~1 wins; $-5$ means Player~2.
Diagonals are checked via 5$\times$5 sub-blocks and \texttt{np.diagonal}.

\subsection{Othello (\texttt{othello.py})}
Valid moves are found by scanning all eight compass directions from each
empty cell.  Disc flipping iterates along each direction until a friendly
disc is found.  The win condition counts discs using
\texttt{np.count\_nonzero} (no loops).

\subsection{Connect Four (\texttt{connect4.py})}
Coins drop to the lowest empty row (found with \texttt{np.where}).
Win detection creates a boolean mask per player and sums all 4-cell
windows using NumPy slicing in four directions.  A sum of~4 on the mask
indicates four-in-a-row.

\subsection{Authentication (\texttt{main.sh})}
Passwords are hashed with \texttt{sha256sum} (or \texttt{shasum -a 256}
on macOS) and stored in \texttt{users.tsv}.  On login the input hash is
compared against the stored hash.  New users are offered registration.

\subsection{Leaderboard (\texttt{leaderboard.sh})}
Reads \texttt{history.csv}, extracts unique players and games using
\texttt{awk}, then loops over each player to count wins and losses per
game.  Results are sorted by the user-selected metric and printed in a
formatted table.

% ─────────────────────────────────────────────────────────────────────
\section{Hurdles Faced and Solutions}

\begin{enumerate}
    \item \textbf{NumPy win detection without loops.}
          Tic-Tac-Toe's diagonal checks were tricky.  We solved it by
          using \texttt{sliding\_window\_view} over 5$\times$5 sub-blocks
          and extracting diagonals with \texttt{np.diagonal}.

    \item \textbf{Connect Four sum ambiguity.}
          Storing players as 1 and 2 meant two player-2 coins could sum
          to~4, falsely matching player~1's target.  We fixed this by
          creating boolean masks per player before summing.

    \item \textbf{Pygame and Matplotlib conflict.}
          Both libraries use display backends.  We used Matplotlib's
          \texttt{Agg} backend to render charts to image files, then
          loaded them into Pygame as surfaces.

    \item \textbf{Cross-platform SHA-256.}
          Linux has \texttt{sha256sum}, macOS has \texttt{shasum -a 256}.
          We detect the available command at runtime.
\end{enumerate}

% ─────────────────────────────────────────────────────────────────────
\section{What We Would Do with 10 More Hours}

\begin{itemize}
    \item Add an AI opponent using minimax with alpha-beta pruning.
    \item Implement network multiplayer so players can compete remotely.
    \item Add sound effects and background music using Pygame's mixer.
    \item Create animated transitions between screens.
    \item Add more games (e.g.\ Checkers, Mancala).
    \item Implement user profiles with avatars and match history.
\end{itemize}

% ─────────────────────────────────────────────────────────────────────
\section{Bibliography}

\begin{enumerate}
    \item Pygame CE Documentation: \url{https://pyga.me}
    \item NumPy Documentation: \url{https://numpy.org/doc/}
    \item Python \texttt{abc} module: \url{https://docs.python.org/3/library/abc.html}
    \item Python Inheritance Tutorial: \url{https://docs.python.org/3/tutorial/classes.html}
    \item Othello Rules: \url{https://www.worldothello.org}
    \item Connect Four Rules: \url{https://boardgamegeek.com/boardgame/2719/connect-four}
    \item SHA-256 in Bash: \texttt{man sha256sum}
    \item Matplotlib Documentation: \url{https://matplotlib.org/stable/}
\end{enumerate}

\end{document}
